\section*{Algebra of Summations}
\addcontentsline{toc}{section}{Algebra of Summations}

The following rules describe the basic algebraic manipulations of summations.
These identities allow sums to be combined, split, and reindexed while
preserving their value.

\begin{center}
\renewcommand{\arraystretch}{1.6}
\begin{tabular}{|l|c|}
\hline
\textbf{Algebra Rule} & \textbf{Identity} \\
\hline

Linearity &
\[
\sum_{k=a}^{b} (a f(k) + b g(k))
=
a \sum_{k=a}^{b} f(k)
+
b \sum_{k=a}^{b} g(k)
\]
\\
\hline

Splitting &
\[
\sum_{k=a}^{c} f(k)
=
\sum_{k=a}^{b} f(k)
+
\sum_{k=b+1}^{c} f(k)
\]
\\
\hline

Index Rename &
\[
\sum_{k=a}^{b} f(k)
=
\sum_{i=a}^{b} f(i)
\]
\\
\hline

Index Shift &
\[
\sum_{k=a}^{b} f(k)
=
\sum_{k=a+1}^{b+1} f(k-1)
\]
\\
\hline

Constant Sum &
\[
\sum_{k=a}^{b} c
=
(b-a+1)c
\]
\\
\hline

Empty Sum &
\[
\sum_{k=a}^{b} f(k) = 0
\qquad \text{if } a>b
\]
\\
\hline

Nested Sums &
\[
\sum_{i=1}^{n}\sum_{j=1}^{m} a_{ij}
=
\sum_{j=1}^{m}\sum_{i=1}^{n} a_{ij}
\quad (\text{when the sums are independent})
\]
\\
\hline

\end{tabular}
\end{center}

\begin{remark}
These rules form the basic algebraic toolkit for manipulating summations.
They are used throughout discrete mathematics, combinatorics, and discrete
calculus, particularly when working with telescoping sums, finite
differences, and summation formulas.
\end{remark}




\section{Algebra of Summations}

The following rules describe the basic algebraic manipulations of summations.
They allow sums to be combined, split, and reindexed while preserving their value.
These identities form the fundamental algebraic toolkit used throughout
discrete mathematics and discrete calculus.

\subsection*{Linearity}

Summation is a linear operator.  For functions \(f\) and \(g\) and constants
\(a\) and \(b\),

\[
\sum_{k=a}^{b} (a f(k) + b g(k))
=
a \sum_{k=a}^{b} f(k)
+
b \sum_{k=a}^{b} g(k).
\]

\begin{remark}
This property allows constants to be factored out of a sum and sums of
expressions to be separated into multiple sums.
\end{remark}

\begin{example}
\[
\sum_{k=1}^{n} (2k + 3)
=
2\sum_{k=1}^{n} k
+
3\sum_{k=1}^{n} 1.
\]
\end{example}

\subsection*{Splitting a Sum}

A sum over an interval may be divided into smaller pieces:

\[
\sum_{k=a}^{c} f(k)
=
\sum_{k=a}^{b} f(k)
+
\sum_{k=b+1}^{c} f(k)
\]

whenever \(a \le b < c\).

\begin{remark}
This rule allows a long summation interval to be broken into parts.
It is frequently used when working with telescoping sums.
\end{remark}

\begin{example}
\[
\sum_{k=1}^{10} f(k)
=
\sum_{k=1}^{5} f(k)
+
\sum_{k=6}^{10} f(k).
\]
\end{example}

\subsection*{Index Renaming}

The index of summation is a \emph{dummy variable}. Its name does not affect
the value of the sum.

\[
\sum_{k=a}^{b} f(k)
=
\sum_{i=a}^{b} f(i).
\]

\begin{remark}
This rule is useful when manipulating expressions containing several sums,
since the index variable can be renamed to avoid conflicts.
\end{remark}

\subsection*{Index Shifting}

It is often convenient to shift the index of summation.

\[
\sum_{k=a}^{b} f(k)
=
\sum_{k=a+1}^{b+1} f(k-1).
\]

\begin{remark}
Index shifting allows sums to be aligned so that terms cancel or combine
properly. This technique appears frequently in proofs involving
recurrence relations and discrete calculus identities.
\end{remark}

\begin{example}
\[
\sum_{k=0}^{n} f(k+1)
=
\sum_{k=1}^{n+1} f(k).
\]
\end{example}

\subsection*{Constant Sum}

If the summand does not depend on the index, the sum simply counts the
number of terms:

\[
\sum_{k=a}^{b} c
=
(b-a+1)c.
\]

\begin{remark}
The quantity \(b-a+1\) represents the number of integers between
\(a\) and \(b\), inclusive.
\end{remark}

\begin{example}
\[
\sum_{k=3}^{7} 5 = 5+5+5+5+5 = 5(5) = 25.
\]
\end{example}

\subsection*{Empty Sum}

If the upper limit is smaller than the lower limit, the sum contains no
terms. By convention,

\[
\sum_{k=a}^{b} f(k) = 0
\qquad \text{whenever } a>b.
\]

\begin{remark}
This convention simplifies many algebraic identities and allows formulas
to remain valid without requiring special cases.
\end{remark}

\subsection*{Nested Sums}

Summations may be nested in a manner similar to nested loops in
programming:

\[
\sum_{i=1}^{n}\sum_{j=1}^{m} a_{ij}.
\]

The inner sum is evaluated first:

\[
\sum_{i=1}^{n} \left(a_{i1}+a_{i2}+\cdots+a_{im}\right).
\]

When the limits of summation are independent, the order of summation may
be exchanged:

\[
\sum_{i=1}^{n}\sum_{j=1}^{m} a_{ij}
=
\sum_{j=1}^{m}\sum_{i=1}^{n} a_{ij}.
\]

\begin{remark}
This property is analogous to exchanging the order of integration in
double integrals.
\end{remark}







\section*{Summation Manipulation Flow Chart}

\begin{center}
\begin{tikzpicture}[
node distance=3cm,
box/.style={draw, rounded corners, align=center, minimum width=3cm, minimum height=1cm},
arrow/.style={->, thick}
]

\node[box] (sum) {$\displaystyle \sum_{k=a}^{b} f(k)$};

\node[box, above=of sum] (linear) {Linearity\\
$\sum (af+bg)=a\sum f+b\sum g$};

\node[box, left=of sum] (split) {Split Interval\\
$\sum_a^c = \sum_a^b + \sum_{b+1}^c$};

\node[box, right=of sum] (shift) {Index Shift\\
$\sum f(k)=\sum f(k-1)$};

\node[box, below=of sum] (rename) {Rename Index\\
$\sum_k f(k)=\sum_i f(i)$};

\node[box, below left=of sum] (constant) {Constant Sum\\
$\sum c=(b-a+1)c$};

\node[box, below right=of sum] (nested) {Nested Sums\\
$\sum_i\sum_j a_{ij}=\sum_j\sum_i a_{ij}$};

\draw[arrow] (sum) -- (linear);
\draw[arrow] (sum) -- (split);
\draw[arrow] (sum) -- (shift);
\draw[arrow] (sum) -- (rename);
\draw[arrow] (sum) -- (constant);
\draw[arrow] (sum) -- (nested);

\end{tikzpicture}
\end{center}

\begin{remark}
Each arrow represents a transformation that preserves the value of the
summation.  These manipulation rules form the basic algebra of summations
and are used constantly when simplifying expressions and proving
identities in discrete mathematics and discrete calculus.
\end{remark}




















\section{Common Summation Identities}

Many sums that appear in discrete mathematics and discrete calculus have
closed-form expressions.  The following identities are among the most
frequently used and are worth memorizing.

\begin{center}
\renewcommand{\arraystretch}{1.7}
\begin{tabular}{|l|c|}
\hline
\textbf{Summation} & \textbf{Closed Form} \\
\hline

Constant Sum &
\[
\sum_{k=1}^{n} 1
=
n
\]
\\
\hline

Sum of First Integers &
\[
\sum_{k=1}^{n} k
=
\frac{n(n+1)}{2}
\]
\\
\hline

Sum of Squares &
\[
\sum_{k=1}^{n} k^2
=
\frac{n(n+1)(2n+1)}{6}
\]
\\
\hline

Sum of Cubes &
\[
\sum_{k=1}^{n} k^3
=
\left(\frac{n(n+1)}{2}\right)^2
\]
\\
\hline

Arithmetic Series &
\[
\sum_{k=0}^{n} (a + kd)
=
(n+1)a + \frac{dn(n+1)}{2}
\]
\\
\hline

Geometric Series &
\[
\sum_{k=0}^{n} r^k
=
\frac{1-r^{n+1}}{1-r}
\quad (r\neq 1)
\]
\\
\hline

Weighted Geometric Series &
\[
\sum_{k=0}^{n} k r^k
=
\frac{r\bigl(1-(n+1)r^n+nr^{n+1}\bigr)}{(1-r)^2}
\quad (r\neq 1)
\]
\\
\hline

Binomial Identity &
\[
\sum_{k=0}^{n} \binom{n}{k}
=
2^n
\]
\\
\hline

\end{tabular}
\end{center}

\begin{remark}
These identities appear frequently in discrete mathematics, combinatorics,
and discrete calculus.  In particular, polynomial summation formulas play
an important role when working with finite differences and Newton series,
while geometric sums arise naturally in recurrence relations and generating
functions.
\end{remark}